% Nouveau EGA V (cote 25): archive cote 25, site page 151 (PDF page 152).
% typescript. Draft compared: none.
% AI-transcribed by Claude Opus 5.5 (claude-opus-5-5, adaptive thinking, effort high), Message Batches API, batch ega25b (msgbatch_01A7n5YuXubUPNK1zsepi8RG), 2026-11-13; original AI transcription, subsequently checked against both supplied scan parts; see PAGES_CHECKED.tsv and CHECK_DECISIONS.tsv.
\scanpage{25/151}
% NOTE: Two separate long diagonal strokes cross the page: one from the upper-right homomorphisme region through the end of the proposition proof, another from the Corollaire 10 heading through the lower-left foot. Their cancellation or working-mark status is unresolved; each block is accompanied by an uncertain diagonal marginal sign.
\mspage{\hand{12}}
\margin{\unsure{$\diagup$}}

projectif de rang m,  $u{:}E \hand{\to} I$ un homomorphisme, d'où un homomorphisme
\[ u\otimes id_{A/I} : E\otimes A/I \hand{\longrightarrow} I/I^2 \ . \]
\xout{Conditions Posons} \inserted[typed]{Soit} $B = A/u(E)A$ , muni de l'augmentation $B\hand{\to} B/J = A/I$
(où $J = I/u(E)A$ ). Les deux conditions suivantes sont équivalentes:

\margin{$\|$} (i) L'anneau \xout{augmenté} B \xout{B/I est} a une augmentation $B\hand{\to} B/J$ \inserted{de type} quadra-%
tique ordinaire, de rang $n-(m-1)= n-m+1$ (le \hand{``}\uuline{rang de l'augmentation}\hand{''}
est \inserted{par définition} celui du Module projectif $J/J^2$ ).

\margin{$|$} (ii) $K = \mathrm{Ker}\ u\otimes id_{A/I}$ est un Module inversible, et l'homomor
\uline{phisme naturel} $K \hand{\longrightarrow} \mathrm{Sym}^2(J/J^2) = \mathrm{Sym}^2( (I/I^2)/\mathrm{Im}\ E\otimes A/I )$ est
``non dégénéré'' .

Il serait peut-être plus clair d'introduire une notation $V = I/I^2$ ,
et $W = J/J^2 = V / \mathrm{Im}\ E\otimes A/I$ . L'homomorphisme qualifié de naturel
dans (ii) s'obtient en notant que K s'envoie en tous cas dans
\margin{$\|$} $I^2/(I^3 +$ \xout{Ixm(E)} $u(IE)) = (I^2/I^3)/\mathrm{Im}\ IE \hand{\simeq} \mathrm{Sym}^2(I/I^2)/\ \mathrm{Im}\ (I/I^2)\otimes(E\otimes A/$
% NOTE: the end of this line runs off the right edge of the scan
où $I/I^2 \otimes(E \otimes A/I) \hand{\to} \mathrm{Sym}^2(I/I^2)$ est défini via l'homomorphisme
canonique de multiplication $I/I^2 \otimes I/I^2 \hand{\to} \mathrm{Sym}^2(I/I^2 )$ . Ces explici-%
tations posées, la démonstration de la prop.9 est à peu près triviale.

\margin{\unsure{$\diagup$}}
\uline{Corollaire} 10. Soit A un anneau local régulier de dimension n, $\underline{m}$ son
idéal maximal, \inserted[typed]{$k=A/\underline{m}$,} E un A-module libre de rang m, $u{:}E \hand{\to} \underline{M}$ un homomorphisme,
\xout{Conditions équivalentes:} $B = A/u(E)A$, $\underline{n} = \underline{m}/u(E)A$ l'idéal maximal de B,
$V = \underline{m}/\underline{m}^2$ , $W = \underline{n}/\underline{n}^2 =$ \xout{km/} $V/\ \mathrm{Im}\ E\otimes k$ . Conditions équivalentes:

\margin{?\}} (i) B est de dimension n-m (i.e. , en termes d'une base \xout{de} $(\hand{\phi}_i)$ de
\margin{$\|$} E, les\xout{t} images des $\hand{\phi}_i$ dans A forment une suite \tikz[baseline=(t.base)]{\node[draw,ellipse,handwritten,text=black,inner sep=1pt] (t) {A-régulière};} et \xout{est} le
\xout{singularité qu} point fermé de Spec(B) est une \uline{singularité quadratique}
ordinaire (par quoi nous entendons que l'augmentation $B\hand{\to} B/\underline{n}$ est \inserted{de type} quadra-%
tique ordinaire).

(ii) Il existe une base $\hand{\phi}_1 ,\ldots., \hand{\phi}_m$ de E telle que \inserted{a)}$A/$\hand{$\sum_{i=1}^{m-1}$}$u(\hand{\phi} i)A$

\mspage{\hand{151}}
